\chapter{Background}

\section{Lean}

A lot of background in Lean was required to undertake this project.
The Lean development exercises the full expressivity of inductive types in
Lean, for example and a thorough understanding of propositional vs definitional
equality which is related to Proof irrelevance [Theorem proving for Lean4] and how that's implemented
in Lean's type theory by special treatment of the lowest Universe [Mario's thesis]

Various representations of formal languages particularly the problem of how
to represent variables in contexts, with strong guarantees from the types System.
[Cite chlipala, HOAS, and lean-lang.org]

\section{Probabilistic Program Semantics}
This section is entirely a summary of the denotational semantics given to Bayesian Probabilistic programs
given by Staton. The probabilistic programs we aim to provide verification tools for in this work
do not go beyond the generality of the semantics presented by Staton. A brief overview of key ideas
is provided below (but it is better to read the )
\begin{itemize}
\item A clear treatment of Staton's commutative semantics.
\item The distinction and intuition behind the difference
between s-finite and sigma-finite measures. Perhaps
even why they're used.
\end{itemize}


\section{Separation Logic}

O'Hearn's original paper. PSL, Lilac, BaSL

It is unclear to me why an axiomatic semantics would need to be provided in terms of measurable maps
from the Hilbert cube. I wonder whether measure theory is not abstract enough and perhaps synthetic
measure theory may hold some intuitions/inspiration or simplifications 
\section{Measure Theory}

\subsection{Pi-lambda theorem}

\subsection{Disintegration theorem}

\section{Symbolic Execution tools}

\section{Iris and MoSeL}

